% Edit the title, individual panel descriptions, and shared notes here.
% All figure placements read these commands; edit once to update every continuation.
\newcommand{\FigSevenTitle}{Scale-resolved risk flux, bottleneck composition, and frozen-model perturbations.}
\newcommand{\FigSevenA}{Distributions of log10 positive risk flux across the 167,633, 63,673, 117,408, and 72,763 directed edges in Scales 1--4, respectively (421,477 edges in total). The vertical axis is probability density rather than an edge count.}
\newcommand{\FigSevenB}{Kernel-density summaries of normalized source pseudotime for top-1\%-flux edges within each scale. Each selected edge contributes its source pseudotime, so sources with several selected outgoing edges are repeated; these are edge-weighted distributions rather than distributions over unique cells. Curves are scaled separately to unit peak height. Pseudotime is dimensionless.}
\newcommand{\FigSevenC}{Fraction of high-pseudotime cells among the 43 unique 24~h sources and 380 unique 48~h targets of top-1\%-flux Scale 3 edges. High pseudotime is defined by the 75th percentile of all cells at the corresponding sampling time. The dashed line marks the approximately 25\% background expectation. Source enrichment is tested using a one-sided Fisher exact test (alternative: greater; odds ratio = 11.2, p = 4.5 \ensuremath{\times} 10\textsuperscript{-13}). The four-asterisk annotation denotes p < 0.0001. The target fraction is a descriptive depletion relative to its 25.1\% background.}
\newcommand{\FigSevenD}{Total absolute change in predicted TrajRisk score across all 6,582 Walesky cells after removing every edge in one scale and rerunning the complete frozen model. The inset summarizes mean absolute TrajRisk-score change by sampling time. Node features and model weights are retained, and no retraining is performed.}
\newcommand{\FigSevenE}{Edge-wise \ensuremath{L_2} distances between source and target 64-dimensional scale embeddings for top-1\%-flux edges and all remaining edges within the same scale. Gray violins represent background edges; colored violins represent top-flux edges. White circles show medians and thick black segments show interquartile ranges. One-sided Mann--Whitney U tests compare top-flux shifts with background shifts (alternative: greater). The Scale 3 annotation (****) denotes p < 0.0001 (actual p = 4.87 \ensuremath{\times} 10\textsuperscript{-173}); n.s. denotes the non-significant Scale 4 result (p = 0.995). Edge counts and all test values are given in \TableRef{S6G}.}
\newcommand{\FigSevenF}{Selected Kyoto Encyclopedia of Genes and Genomes (KEGG) pathway enrichments from separate up- and downregulated gene sets in bottleneck versus non-bottleneck target cells at the same sampling time. Upward and downward triangles denote up- and downregulated gene sets, respectively; the upper section contains upregulated sets and the lower section downregulated sets. Marker area increases with \ensuremath{-}log10 of the Enrichr-adjusted p value and is clipped to a display range. Colors identify temporal scales. The displayed pathway subset is not a complete enrichment report.}
\newcommand{\FigSevenNotes}{Scales 1--4 connect 0--6, 6--24, 24--48, and 48--96~h, respectively. Risk flux and risk amplification summarize the model-derived trajectory-state score; scores and their absolute changes are dimensionless. Source and target counts are unique cells; edge-level comparisons are exploratory summaries of the trained model. A bottleneck identifies a model-derived candidate transition and does not establish a cell's fate or a causal recovery mechanism.}
